\documentclass[../../main.tex]{subfiles}
\begin{document}

{\bf [解答]}

与えられた条件
\begin{align}
  \begin{cases}
    p, q \in \mathbb{N} \\
    -q \le b \le 0 \le a \le p \\
    b \le c \le a \\
    a, b, c \in \mathbb{Z}
  \end{cases} \label{eq:1}
\end{align}
について考える．

\paragraph{(1)}

% TODO: (a,b)平面で\cref{eq:1}を図示，一旦cは無視
\begin{figure}[htb]
  \centering
  \begin{tikzpicture}[x=2.5cm, y=2.0cm, >=stealth, every node/.style={font=\small}]

    % ============================================================
    % 1. 定数設定 (例: p = 1.6, q = 1.2)
    % ============================================================
    \def\pVal{1.6}
    \def\qVal{1.2}

    \coordinate (O) at (0, 0);
    \coordinate (Amax) at (\pVal, 0);
    \coordinate (Bmin) at (0, -\qVal);
    \coordinate (Corner) at (\pVal, -\qVal);

    % ============================================================
    % 2. 領域 (0 <= a <= p, -q <= b <= 0) の斜線ハッチング
    % ============================================================
    \fill[pattern=north east lines, pattern color=blue!70, opacity=0.7]
    (O) -- (Amax) -- (Corner) -- (Bmin) -- cycle;

    % ============================================================
    % 3. 座標軸
    % ============================================================
    \draw[->, thick, gray!80!black] (-0.3, 0) -- (\pVal + 0.5, 0) node[right] {$a$};
    \draw[->, thick, gray!80!black] (0, -\qVal - 0.4) -- (0, 0.6) node[above] {$b$};
    \node[above left=2pt] at (O) {$O$};

    % ============================================================
    % 4. 領域の境界線 (長方形)
    % ============================================================
    \draw[line width=1.4pt, blue!85!black] (O) -- (Amax) -- (Corner) -- (Bmin) -- cycle;

    % 目盛りラベル
    \node[above=2pt, font=\scriptsize] at (Amax) {$p$};
    \node[left=2pt, font=\scriptsize] at (Bmin) {$-q$};

    % ============================================================
    % 5. (1) の直線 a + b = p および a + b = -q
    % ============================================================
    % 直線 a + b = p (点 (p, 0) を通り傾き -1)
    \draw[dashed, red!85!black, line width=1.1pt]
    ($ (Amax) + (-0.5, 0.5) $) -- ($ (Amax) + (0.3, -0.3) $)
    node[right, font=\scriptsize\bfseries] {$a+b = p$};
    \node[circle, fill=red!85!black, inner sep=1.5pt] at (Amax) {}; \node[below right=2pt, font=\scriptsize\bfseries, red!90!black] at (Amax) {$(p, 0)$};

    % 直線 a + b = -q (点 (0, -q) を通り傾き -1)
    \draw[dashed, teal!85!black, line width=1.1pt]
    ($ (Bmin) + (-0.3, 0.3) $) -- ($ (Bmin) + (0.5, -0.5) $)
    node[below right, font=\scriptsize\bfseries] {$a+b = -q$};
    \node[circle, fill=teal!85!black, inner sep=1.5pt] at (Bmin) {}; \node[above left=2pt, font=\scriptsize\bfseries, teal!90!black] at (Bmin) {$(0, -q)$};

  \end{tikzpicture}
  \caption{$(a, b)$ 平面における基本領域（青斜線部）と極値をとる点 $(p, 0), (0, -q)$}
  \label{fig:ab_plane_basic_domain}
\end{figure}

\subparagraph{$w([a, b; c])=-q$について}

まず前者の条件についてかんがえる．
与えられた条件は
\begin{align}
  w([a, b; c]) = p - q - (a+b) = -q \quad \therefore p = (a+b) \label{eq:2}
\end{align}
である．\cref{eq:1}から$b \le 0, a \le p$だから，\cref{eq:2}を満たすには
$(a, b) = (p, 0)$ が必要で，\cref{eq:1}から$c = 0, 1, \dots, p$ とすれば十分である．
よって，求める個数は$p+1$ 個である．

\subparagraph{$w([a, b; c])=p$について}

次に後者の条件についてかんがえる．
与えられた条件は
\begin{align}
  W([a, b; c]) = p - q - (a+b) = p \quad \therefore -q = (a+b) \label{eq:3}
\end{align}
である．\cref{eq:1}から $-q \le b, \ 0 \le a$ だから，\cref{eq:3}を満たすには
$(a, b) = (0, -q)$ が必要で，\cref{eq:1}から $c = 0, -1, \dots, -q$ とすれば十分である．
よって求める個数は$q+1$ 個である．
\casesend

以上から
\begin{align}
  \begin{cases}
    w([a, b; c])=-q: p+1\text{個} \\
    w([a, b; c])=p: q+1\text{個}
  \end{cases}
\end{align}
である．

\paragraph{(2)}

$p = q$ の時，$s \in \mathbb{Z}$ に対して
\begin{align}
  W([a, b; c]) = -(a+b) = -p + s \quad \therefore a+b = p-s \label{eq:4}
\end{align}
である．\cref{eq:1}より $-p \le a+b \le p$ だから，
\begin{align}
  -p \le p-s \le p \quad \therefore 0 \le s \le 2p \label{eq:5}
\end{align}
である．\cref{eq:4}から$b = p-s-a$ だから\cref{eq:1}に代入して
\begin{align}
  -p \le p-s-a \le 0 \iff p-s \le a \le 2p-s \label{eq:6}
\end{align}
の時のみ，各 $a$ に対して $b$ がただ1つ対応する．
\cref{eq:1}より
\begin{align}
  0 \le a \le p \label{eq:7}
\end{align}
でもあるから，\cref{eq:6,eq:7}を満たす$a$について，
\cref{eq:5}の範囲で$s$を動かして検討する．

% TODO: \cref{eq:6,eq:7}について$s$の大きさで場合分け（sa平面に図示）
% TODO: \cref{eq:6,eq:7}について$s$の大きさで場合分け（sa平面に図示）
\begin{figure}[htb]
  \centering
  \begin{tikzpicture}[x=2.4cm, y=2.4cm, >=stealth, every node/.style={font=\small}]

    % ============================================================
    % 1. 座標定義 (p = 1.0 とする)
    % ============================================================
    \def\pVal{1.0}

    \coordinate (O)       at (0, 0);
    \coordinate (P_s)     at (\pVal, 0);       % (p, 0)
    \coordinate (TwoP_s)  at ({2*\pVal}, 0);   % (2p, 0)
    \coordinate (P_a)     at (0, \pVal);       % (0, p)
    \coordinate (TwoP_a)  at (0, {2*\pVal});   % (0, 2p)
    \coordinate (Mid)     at (\pVal, \pVal);   % (p, p)

    % ============================================================
    % 2. 満たすべき領域 (斜線ハッチング: スケッチ通り)
    % 0 <= s <= p では: p - s <= a <= p
    % p <= s <= 2p では: 0 <= a <= 2p - s
    % ============================================================
    \fill[pattern=north east lines, pattern color=blue!75, opacity=0.75]
    (P_a) -- (Mid) -- (P_s) -- cycle;
    \fill[pattern=north east lines, pattern color=blue!75, opacity=0.75]
    (Mid) -- (TwoP_s) -- (P_s) -- cycle;

    % ============================================================
    % 3. 座標軸
    % ============================================================
    \draw[->, thick, gray!80!black] (-0.2, 0) -- ({2*\pVal + 0.4}, 0) node[below] {$s$};
    \draw[->, thick, gray!80!black] (0, -0.2) -- (0, {2*\pVal + 0.3}) node[left] {$a$};
    \node[below left=2pt] at (O) {$0$};

    % ============================================================
    % 4. 境界線と直線群（スケッチ通り）
    % ============================================================
    % a = 2p - s (点線〜実線)
    \draw[dashed, gray!60, line width=0.9pt] (TwoP_a) -- (Mid);
    \draw[line width=1.3pt, blue!85!black] (Mid) -- (TwoP_s)
    node[midway, above right=1pt, font=\footnotesize\bfseries] {$a = 2p - s$};

    % a = p - s (実線)
    \draw[line width=1.3pt, blue!85!black] (P_a) -- (P_s);

    % a = p (水平線: (0, p) から (p, p) まで実線、その先は破線)
    \draw[line width=1.3pt, blue!85!black] (P_a) -- (Mid);
    \draw[dashed, gray!50] (Mid) -- ({2*\pVal + 0.1}, \pVal);

    % 場合分けの境界線 s = p (縦の破線: (p, 0) から (p, p))
    \draw[dashed, red!85!black, line width=1.1pt] (P_s) -- (Mid);

    % ============================================================
    % 5. 軸の目盛り・点プロット
    % ============================================================
    \node[below=3pt, font=\normalsize] at (P_s) {$p$};
    \node[below=3pt, font=\normalsize] at (TwoP_s) {$2p$};
    \node[left=3pt, font=\normalsize]  at (P_a) {$p$};

    % 式のラベル（手書きスケッチの引き出し線風）
    \node[below=14pt, font=\footnotesize\bfseries] at (0.5, 0.45) {$a = p - s$};

    % 場合分けの領域ラベル (i), (ii)
    \node[font=\footnotesize\bfseries, blue!90!black] at (0.4, 0.75) {(i)};
    \node[font=\footnotesize\bfseries, blue!90!black] at (1.2, 0.4)  {(ii)};

  \end{tikzpicture}
  \caption{$(s, a)$ 平面における条件の領域（斜線部）と $s = p$ による場合分け}
  \label{fig:sa_plane_cases_sketch}
\end{figure}

\subparagraph{(i) $0 \le s \le p$ の時}

この時は\cref{eq:6,eq:7}は
\begin{align}
  p-s \le a \le p
\end{align}
となる．各 $(a,b)$ に対し，$c$ は
\begin{align}
  a - b + 1 = a - (p-s-a) + 1 = 2a + (1+s-p) \text{ 通り}
\end{align}
あるから，全部で
\begin{align}
  \sum_{a=p-s}^{p} \{ 2a + (1+s-p) \}
  &= \dfrac{1}{2} \cdot 2 \cdot (2p-s)(s+1) + (1+s-p)(s+1) \\
  &= (p+1)(s+1) \label{eq:8}
\end{align}
通りである．

\subparagraph{(ii) $p \le s \le 2p$ の時}

この時は\cref{eq:6,eq:7}は
\begin{align}
  0 \le a \le 2p-s
\end{align}
となる．(i)と同様に，全ての場合の数は
\begin{align}
  \sum_{a=0}^{2p-s} \{ 2a + (1+s-p) \}
  &= \dfrac{1}{2} \cdot 2 \cdot (2p-s)(2p-s+1) + (1+s-p)(2p-s+1) \\
  &= (p+1)(2p-s+1) \label{eq:9}
\end{align}
である．
\casesend

以上\cref{eq:8,eq:9}をまとめて，求める数は
\begin{align}
  \begin{cases}
    (p+1)(s+1) & (0 \le s \le p) \\
    (p+1)(2p-s+1) & (p \le s \le 2p) \\
    0 & (\text{otherwise})
  \end{cases} \label{eq:10}
\end{align}
である．

\paragraph{(3)}

\cref{eq:10}を$s$ について和をとれば良く，$p \ge 1$ から
\begin{align}
  &\sum_{s=0}^{p-1} (p+1)(s+1) + \sum_{s=p}^{2p} (p+1)(2p-s+1) \\
  &= (p+1) \sum_{s=1}^{p} s + (p+1) \sum_{s=1}^{p+1} s \\
  &= (p+1) \left[ \dfrac{1}{2}p(p+1) + \dfrac{1}{2}(p+1)(p+2) \right] = (p+1)^3
\end{align}
である．

\end{document}
